\documentclass[10pt,a4paper]{amsart}
\usepackage[margin=19mm]{geometry}
\usepackage{amsmath,amssymb,mathtools}
\usepackage[scr=rsfs,cal=euler]{mathalfa}
\usepackage{xfrac}
\usepackage[unicode,hidelinks,bookmarks=false]{hyperref}
\newcommand{\ie}{i.e.\ }
\newcommand{\cf}{cf.\ }
\newcommand{\resp}{respectively\ }
\newcommand{\Subsection}{\subsection}
\providecommand{\nobreakdash}{\nobreak\hbox{-}}
\setlength{\emergencystretch}{2em}
\multlinegap0pt
\usepackage{bm}
\newcommand{\R}{\mathbb{R}}
\newcommand{\norm}[1]{\lVert#1\rVert}
\newcommand{\bUnk}{\mathcal{U}_{n,k}}
\newcommand{\la}{\big\langle}
\newcommand{\ra}{\big\rangle}
\newcommand{\bmF}{\bm{F}}
\newcommand{\bmu}{\bm{\mu}}
\newcommand{\cA}{\mathcal{A}}
\theoremstyle{plain}
\newtheorem{fact}{Fact}
\newtheorem{theorem}{Theorem}
\newtheorem{lemma}[theorem]{Lemma}
\newtheorem{corollary}[theorem]{Corollary}
\newtheorem{proposition}[theorem]{Proposition}
\newtheorem{remark}{Remark}
\theoremstyle{definition}
\newtheorem{definition}[theorem]{Definition}
\newtheorem{example}[theorem]{Example}
\newtheorem{experiment}[theorem]{Experiment}
\title{BalLOT: Balanced $k$-means clustering with optimal transport}
\author{Wenyan Luo and Dustin G. Mixon}
\date{Source: arXiv:2512.05926v2 (CC BY 4.0); the authors' own native \LaTeX{} source, set here in the editorial AMS \texttt{amsart} wrapper (the source's own class is \texttt{article}, which is not redistributed).}
\begin{document}
\maketitle
\noindent\textit{Editorial scope.} The counted claim of this unit is the article's Theorem 14 (source label \texttt{thm.generic initialization by partition}), the unique-minimizer statement for the first assignment step of balanced $k$-means with optimal transport when the initial centroids are the centroids of a balanced partition of generic data, delivered together with the authors' complete original proof. The counted proof is delivered as the single primary proof body, and its statement and proof are the only retained passages: every one of their characters is the authors' own. Printed numbering: the authors' own theorem-environment declarations are carried unchanged, including the shared counter of the theorem, lemma, corollary, proposition, definition, example and experiment environments, and the counter is advanced so that the delivered PDF prints the counted claim as Theorem~14, exactly as the published article does. Omitted: the rest of the article, that is the abstract, the introduction, the model and algorithm sections, the earlier Theorem~13 with its proof, the numerical experiments and the appendices. Theorem~13 is omitted even though the first sentence of the counted proof refers to it in plain words, because the retained lines of the passage around Theorem~13 contain two cross-references to results that lie in the appendix of the article and that cannot resolve inside a source-local unit of this size; the counted proof is nevertheless complete in itself, since it re-derives the genericity identity and the exchange property that it uses. Class substitution: the source class is the standard \texttt{article} class with the authors' own packages; the wrapper is AMS \texttt{amsart} carrying the authors' own macros, including their blackboard-bold macros and the \texttt{bm} package for the bold symbols used in the retained passage. Reference handling: no retained passage contains a citation, the article's bibliography exists only in the authors' file \texttt{references.bib}, which is kept verbatim in the eprint directory and is not redistributed, and no reference entry is transcribed or invented. No mathematics is written, repaired or re-derived here.
\setcounter{theorem}{13}
\begin{theorem}[Initialization by partition of data]\label{thm.generic initialization by partition}
Suppose $\{\bm{x}_i\}_{i\in[n]}$ is generic and $\bmu^0$ consists of the centroids of a balanced partition of the $\bm{x}_i$'s.
Then the minimizer of $f(\bmF,\bmu^0)$ subject to $\bmF\in\bUnk$ is unique.
\end{theorem}

\begin{proof}
Suppose the minimizer is not unique.  As in the previous proof, the optimal face then contains distinct integral vertices $\bm{P}^1,\bm{P}^2\in\bUnk$ for which $f(\bm{P}^1,\bmu^0)=f(\bm{P}^2,\bmu^0)$.  It suffices to show that these couplings are suboptimal for $f(\cdot,\bmu^0)$.
Suppose the partition of $[n]$ that determines $\bmu^0$ is given by 
\[
C_1\sqcup \cdots \sqcup C_k 
= [n].
\]
For $\ell\in \{1,2 \}$, we define $\alpha_\ell\colon [n]\to \{ C_1, C_2,\dotsc, C_k\}$ to be the set-valued assignment by $\bm{P}^\ell$, that is,
\[
\alpha_\ell(i) 
= C_j\quad \text{if $\bm{P}^\ell$ assigns $\bm{x}_i$ to the centroid of $\{\bm{x}_{i'}\}_{i'\in C_j}$}.
\] 
Then our assumption $f(\bm{P}^1,\bmu^0)=f(\bm{P}^2,\bmu^0)$ is equivalent to $p(\bm{X})=0$, where
\[
p(\bm{X}) 
= \sum_{i\in \cA} \la \bm{x}_i , \sum_{j\in \alpha_1(i)} \bm{x}_j - \sum_{j' \in \alpha_2(i) } \bm{x}_{j'}  \ra,
\]
and $\cA := \{ i\in [n]: \alpha_1(i) \neq \alpha_2(i)\}$.
As before, since $p(\bm{X})=0$ and $\bm{X}$ is generic, it follows that $p(\bm{X})$ is identically zero.
In particular, the following identity holds when treating the $\bm{x}_i$'s as formal variables:
\[
\sum_{i\in \cA}\sum_{j\in \alpha_1(i)} \la \bm{x}_i ,  \bm{x}_j \ra 
= \sum_{i\in \cA} \sum_{j' \in \alpha_2(i) } \la \bm{x}_i,\bm{x}_{j'}  \ra.
\]
Notably, every term on the left-hand side corresponds to a term on the right-hand side, and vice versa.
Since $\alpha_1(i)$ and $\alpha_2(i)$ are disjoint whenever $\alpha_1(i)\neq\alpha_2(i)$, it follows that
\[
\text{$i\in \cA$ ~and~ $j\in \alpha_1(i)$}
\quad
\Longleftrightarrow
\quad
\text{$j \in \cA$ ~and~ $i\in \alpha_2(j)$}.
\]
We refer to this equivalence as the \textit{exchange property}.
In what follows, we repeatedly appeal to the exchange property in order to uncover how $\alpha_1$ and $\alpha_2$ behave over $\cA$.
(Throughout, we write $i\sim j$ if $i$ and $j$ belong to the same $C_\ell$.)

First, we claim that $\cA$ is a union of $C_\ell$'s. 
Pick any $i\in\cA$ and any $i'\sim i$.
If we select $j\in\alpha_1(i)$, then by the exchange property (in the forward direction), we have $j\in\cA$ and $i\in\alpha_2(j)$.
Since $i'\sim i$, we also have $i'\in\alpha_2(j)$.
Since $j\in\cA$ and $i'\in\alpha_2(j)$, the exchange property (in the reverse direction) then gives that $i'\in\cA$.

Next, we claim that for each $\ell\in\{1,2\}$, $\alpha_\ell$ maps points in $\cA$ to points in $\cA$.
The $\ell=1$ case follows from the exchange property in the forward direction, while the $\ell=2$ case follows from the exchange property in the reverse direction.

Next, we claim that $i\sim i'$ in $\cA$ implies $\alpha_\ell(i)=\alpha_\ell(i')$ for both $\ell\in\{1,2\}$, i.e., each $\alpha_\ell$ maps clusters in $\cA$ to clusters in $\cA$.
We will prove this for $\ell=1$, as the proof for $\ell=2$ is identical.
For $i\sim i'$ in $\cA$, the exchange property (in the forward direction) gives that for every $j\in\alpha_1(i)$ and $j'\in\alpha_1(i')$, it holds that $j,j'\in\cA$ and $i\in\alpha_2(j)$ and $i'\in\alpha_2(j')$, in which case $i\sim i'$ forces $\alpha_2(j)=\alpha_2(j')$.
So we have $j\in\cA$ and $i,i'\in \alpha_2(j)$.
Then the exchange property (in the reverse direction) gives that $j$ is in both $\alpha_1(i)$ and $\alpha_1(i')$, and so $\alpha_1(i)=\alpha_1(i')$.

Finally, we claim that for each $\ell\in\{1,2\}$, if $i,i'\in \cA$ and $\alpha_\ell(i) = \alpha_\ell(i')$, then $i\sim i'$, i.e., each $\alpha_\ell$ permutes the clusters in $\cA$.
(Again, we only prove this for $\ell=1$.)
Indeed, take any $j\in \alpha_1(i) = \alpha_1(i')$.
By exchange property (in the forward direction), it follows that $j\in \cA$ and $i,i'\in  \alpha_2(j)$, meaning $i\sim i'$.

At this point, we know that $\alpha_1$ and $\alpha_2$ both permute the clusters in $\cA$. 
Next, since $\bm{X}$ is generic, the cluster centroids are distinct, and so equality in the lower bound
\begin{align*}
\sum_{a\in [k]} \sum_{\substack{i\in\cA\\\alpha_1(i) = C_a}} \| \bm{x}_i - \bmu^0_a\|^2 
&= \sum_{a\in [k]}
\sum_{\substack{i\in\cA\\\alpha_1(i) = C_a}} \Bigg(
    \bigg\|\bm{x}_i - \frac{1}{n/k} \sum_{\substack{j\in\cA\\\alpha_1(j) = C_a}} \bm{x}_j \bigg\|^2 + \bigg\|\bmu^0_{a} -\frac{1}{n/k} \sum_{\substack{j\in\cA\\\alpha_1(j) = C_a}} \bm{x}_j \bigg\|^2 \Bigg)
    \\
    & \ge \sum_{a\in [k]} \sum_{\substack{i\in\cA\\\alpha_1(i) = C_a}} \bigg\|\bm{x}_i - \frac{1}{n/k} \sum_{\substack{j\in\cA\\\alpha_1(j) = C_a}} \bm{x}_j \bigg\|^2
\end{align*} 
occurs precisely when every $\bmu^0_a$ is the centroid of points indexed by $C_a$, i.e., $\alpha_1$ maps every cluster in $\cA$ to itself. 
Similarly, the same holds for $\alpha_2$. 
Since $\bm{P}^1\neq \bm{P}^2$, it necessarily holds that $\alpha_1\neq \alpha_2$ on $\cA$, and so we can decrease $f(\cdot, \bmu^0)$ by changing one of them to the identity cluster permutation on $\cA$.
\end{proof}

\end{document}
